\documentclass{article}
\usepackage[utf8]{inputenc}

\title{Computación Distribuida \\ Tarea 03}
\author{Olea Olvera Marco Iván \\ Arteaga Hernández Iabin}
\date{}

\usepackage[top=1.5cm, bottom=1.5cm, left=2cm, right=2cm]{geometry}
\usepackage{fouriernc}
\usepackage{amssymb}
\usepackage{listings}
\lstset{
  basicstyle=\ttfamily,
  mathescape
}

\usepackage{graphicx}
\graphicspath{{./}}

\begin{document}

\maketitle

\section{}

Sea $n$ el número de procesos. \\ \\
\textbf{Terminación.} Solamente hay una ronda y termina cuando cada proceso $p_i$ haya mandado el par $(v_i,\ i)$ a todos sus vecinos. Es decir, $p_i$ decidirá algún valor después de $n - 1$ pasos (porque la gráfica es completa). \\ \\
\textbf{Validez.} El arreglo $V_i$ fue inicializado con un valor no vacío $v_i$, el cual es la propuesta de $p_i$. El arreglo será modificado a lo largo de la ronda pero no se eliminarán valores, sólo se pueden agregar valores propuestos por otros procesos. Esto implica que al final de la única ronda, $V_i$ contendrá al menos un valor que fue propuesto por otro proceso. Por lo tanto, todos los procesos sin fallas decidirán un valor que fue propuesto por algún proceso. \\ \\
\textbf{Acuerdo} Se cumple si no ocurren fallas, o si un proceso falla y envió cero o bien todos sus mensajes. Sea $a \in \{1,2,...,n\}$ el mínimo tal que $p_a$ envió todos sus mensajes. Entonces $V_i[a] = v_a$ para todo proceso correcto $p_i$. Por lo tanto, todos los procesos correctos decidirán $v_a$. \\

\section{}

Podría no funcionar si un proceso $p_a$ envió su mensaje a uno de sus vecinos pero falló antes de enviarlo al resto de sus vecinos. Si existe otro proceso $p_b$ con $b < a$ que sí logró enviar su mensaje a todos sus vecinos, el protocolo funcionará. Sin embargo, si $p_a$ fue el mínimo que envió al menos un mensaje, significa que el protocolo no funcionó porque los procesos que recibieron $(v_a,\ a)$ decidieron $v_a$ y los demás decidieron otra cosa.

\section{}

Si un proceso $p_a$ falla entonces se requiere de otra ronda para que aquellos que hayan recibido el mensaje de $p_a$ puedan reenviarlo a los demás procesos.

Sí puede suceder que las $f$ fallas ocurran en la primera ronda y, por tanto, al terminar la segunda ronda llegarán a un acuerdo, pero esto sería el mejor caso. El peor caso es que ocurra una falla por ronda y consecuentemente la ronda $f + 1$ será la primera sin fallas. 

Por lo tanto, es necesario ejecutar $f + 1$ rondas para garantizar que se llegará a un acuerdo.

\section{}

Recordatorio: Dado un proceso $p_i$, $R_i(r)$ es el conjunto de procesos de los cuales se recibió mensajes durante la ronda $r$. Por definición, $R_i(0) = \Pi$ para todo $i$. \\
Además, $f$ es el número de fallas permitidas mientras que $t$ es el número real de procesos que fallaron. \\

Por hipótesis de esta pregunta, $t = 0$. Lo único que se sabe de $f$ es que $f < n$. Por tanto tenemos dos casos:

\begin{enumerate}
    \item Si $f = 0$, entonces el protocolo termina en $f + 1 = 1$ ronda, según la descripción del algoritmo.
    \item Si $f > 0$, entonces en la ronda $r = 1$ tenemos que $R_i(r) = R_i(1) = \Pi$ para todo $i$, ya que ningún proceso falla ($t = 0$). Pero $\Pi = R_i(0)$, de donde $R_i(r) = R_i(1) = R_i(0) = R_i(r - 1)$. Además como $r = 1 < f + 1$ (porque $f > 0$), entonces se prenderá la bandera que le indicará a $p_i$ que debe decidir al final de la ronda $r + 1 = 2$. \\
    Por lo tanto, el protocolo termina en dos rondas.
\end{enumerate}

\end{document}
